% =====================================================================
%  CARNET D'ENTRAÎNEMENT — CHIMIE — FICHE 3 — THÈME 2
%  Composition : atomes et ions — NIVEAU MOYEN — ÉNONCÉS
% =====================================================================
\documentclass[11pt,a4paper]{article}

\usepackage[utf8]{inputenc}
\usepackage[T1]{fontenc}
\usepackage{lmodern}
\usepackage[french]{babel}

\usepackage{amsmath,amssymb,mathtools}
\usepackage{xcolor}
\usepackage{colortbl}
\usepackage{tikz}
\usetikzlibrary{arrows.meta,positioning,calc}
\usepackage[most]{tcolorbox}
\usepackage{enumitem}
\usepackage{booktabs}
\usepackage{array}
\usepackage[a4paper,margin=2.1cm,headheight=26pt,headsep=9pt]{geometry}
\usepackage{fancyhdr}
\usepackage[hidelinks]{hyperref}

\definecolor{primary}{RGB}{124,78,140}
\definecolor{primarydark}{RGB}{74,44,92}
\definecolor{defcol}{RGB}{0,114,178}
\definecolor{thmcol}{RGB}{0,140,120}
\definecolor{attncol}{RGB}{200,40,55}

\newcommand{\nuc}[3]{\prescript{#1}{#2}{\mathrm{#3}}}
\newcommand{\pp}{p^{+}}
\newcommand{\ee}{e^{-}}
\newcommand{\nz}{n^{0}}

\pagestyle{fancy}
\fancyhf{}
\fancyhead[L]{\small\textcolor{primarydark}{\textbf{Fiche 3 — Thème 2}} : Composition des atomes et des ions}
\fancyhead[R]{\small\textcolor{primarydark}{\textsc{Moyen — Énoncés}}}
\fancyfoot[C]{\small\thepage}
\renewcommand{\headrulewidth}{0.8pt}
\renewcommand{\headrule}{\hbox to\headwidth{\color{primary}\leaders\hrule height \headrulewidth\hfill}}

\newcounter{exo}
\newcommand{\exo}[1][]{\stepcounter{exo}\par\medskip%
  \noindent{\color{primary}\rule[-2pt]{3.5pt}{15pt}}\ \textbf{\color{primarydark}\large Exercice \theexo}%
  \ifx\relax#1\relax\relax\else\ \textmd{\normalsize\color{black!65}— \itshape #1}\fi\par\nopagebreak\smallskip}

\setlist{leftmargin=1.7em,topsep=2pt,itemsep=3pt}

\begin{document}

\begin{tikzpicture}
\node[rectangle,rounded corners=6pt,fill=primary,text=white,
      minimum width=\textwidth,minimum height=1.35cm,align=center,inner sep=6pt]
  {\Large\bfseries Thème 2 — Composition des atomes et des ions\\[2pt]
   \normalsize\mdseries Niveau Moyen — Énoncés};
\end{tikzpicture}

\vspace{3pt}
{\small\itshape Règle d'or : pour un ion de charge $q$ (signée), $\ee = Z - q$, tandis que $\np = Z$ et
$\nz = A - Z$ \textbf{ne changent pas} (le noyau reste identique).}

% ---------------------------------------------------------------
\exo[un cation]
Le cation magnésium $\mathrm{Mg}^{2+}$ provient de l'isotope $\nuc{24}{12}{Mg}$. Donne son nombre de
protons, de neutrons et d'électrons.

% ---------------------------------------------------------------
\exo[un anion]
L'anion oxyde $\mathrm{O}^{2-}$ provient de l'isotope $\nuc{16}{8}{O}$. Donne son nombre de protons, de
neutrons et d'électrons. En quoi le résultat pour les électrons diffère-t-il de celui d'un cation ?

% ---------------------------------------------------------------
\exo[composition d'un ion à partir de sa notation]
Détermine le nombre de protons, de neutrons et d'électrons de l'ion $\nuc{15}{7}{N}^{3-}$.

% ---------------------------------------------------------------
\exo[un ion plus lourd]
On considère l'ion baryum $\nuc{132}{56}{Ba}^{2+}$. Combien possède-t-il de protons, de neutrons et
d'électrons ?

% ---------------------------------------------------------------
\exo[atome contre ion, et le piège du noyau]
On étudie le scandium $\nuc{45}{21}{Sc}$ et son cation $\mathrm{Sc}^{3+}$.
\begin{enumerate}[label=\alph*)]
  \item Combien d'électrons possède l'atome neutre $\nuc{45}{21}{Sc}$ ? Et le cation
        $\mathrm{Sc}^{3+}$ ?
  \item Le nombre de protons et le nombre de neutrons changent-ils entre l'atome et son cation ?
  \item Combien d'électrons se trouvent \textbf{dans le noyau} de $\mathrm{Sc}^{3+}$ ? (Attention, c'est
        un piège classique.)
\end{enumerate}

\end{document}
